\documentclass{article}
\usepackage[english]{babel}
\usepackage{amsmath}

\title{On the Behaviour of Stochastic Processes in Distributed Systems}
\author{J.~Smith}
\date{2026}

\begin{document}
\maketitle

\section{Introduction}

The behaviour of distributed systems has been widely studied in the literature over
the past several decades.  In this paper we analyse the colour of noise arising in
large-scale networks.  The programme of research initiated by earlier authors is
extended here to cover generalised models that incorporate both synchronisation
and randomisation primitives.

% STYLE-002: double space after period
It is well known that consensus is hard.  However,  the problem can be decomposed
into smaller sub-problems that are each tractable in isolation.

% STYLE-017: sentence > 40 words
The fundamental observation underlying our approach is that when the variance of the
delay distribution exceeds the threshold predicted by the central limit theorem applied
to the sum of independent identically distributed random variables drawn from a
heavy-tailed family then the system transitions into an unstable regime characterised
by unbounded message queues.

% STYLE-022: passive voice
The algorithm was implemented in OCaml and the results were verified by running
the test suite on a cluster of 64 machines.  The data was collected over a period
of three weeks and the analysis was performed using standard statistical methods.

\section{Methodology}

We randomize the initial configuration and optimize the objective function using
gradient descent.  The color of the output signal is measured at each step.
% Note: American spellings above — consistent within section

% CE-001/CE-002: Canadian/British mix with American
The behaviour of the optimiser was analysed carefully.  We then maximize the
throughput while minimizing latency, but also favour the centralised approach.

% LANG-014: -ise vs -ize inconsistency
We standardise our notation but also utilize the framework from~\cite{jones2024}.
The organisation recognises the need to modernize its infrastructure.

% LANG-015/LANG-016: programme vs program
The research programme outlined in Section~1 describes a computer program that
automates the scheduling of the programme's milestones.

\section{Results}

Let $f(x) = \sum_{i=1}^{n} x_i^2$ denote the objective.  We minimise $f$ subject
to the constraint $\|x\| \le 1$.  The optimisation converges in $O(n \log n)$ steps.

% Correct pattern: consistent British spelling
The generalisation of Theorem~2 to the multi-dimensional case is straightforward.
The parametrisation of the model was chosen to maximise interpretability.

\section{Conclusion}

In conclusion, the behaviour of distributed stochastic systems can be characterised
by a small number of parameters.  Future work will focus on the parallelisation of
the algorithm and its deployment in containerised environments.

\end{document}
